% Môn: Vật lý
% Lớp: 11
% Chương: Chương I. Dao động
% Bài: L11C1 Bài 1. Dao động điều hòa · Chưa có dạng
% ===== Câu 14 =====
% ID: L11C1B1-181-DS
% Mức: TH
\dangbt{Chưa có dạng}
\begin{ex}
% ID: L11C1B1-181-DS
% Mức: TH
Phương trình dao động điều hoà của một vật là $x=5\cos\left(10\pi t-\dfrac{\pi}{2}\right)\,(\mathrm{cm})$. Tính thời gian để vật đi được quãng đường $2,5\,\mathrm{cm}$ kể từ thời điểm $t=0$. $\nguon{ly11 kntt.pdf}$
\choiceTF

\loigiai{
Chu kì $T=\dfrac{2\pi}{\omega} = \dfrac{1}{5}$ s.

Tại thời điểm ban đầu t = 0 thì $x_{1} = 5\cos\left( {10\pi.0 - \dfrac{\pi}{2}} \right) = 0\left( \text{cm} \right)$

Từ Hình 2.2 $G$, ta thấy thời gian để vật đi được quãng đường $2,5\,\mathrm{cm}$ kể từ t = 0 (từ x₁ = 0 đến x₂ = $2,5\,\mathrm{cm}$) là: t = $\dfrac{\text{T}}{12} = \dfrac{0,2}{12} = \dfrac{1}{60}\text{s}$
}
\end{ex}

% ===== Câu 16 =====
% ID: L11C1B1-182-TN
% Mức: NB
\begin{ex}
% ID: L11C1B1-182-TN
% Mức: NB
Một chất điểm dao động điều hoà có tần số góc $\omega = 10\pi(\mathrm{rad/s})$. Tần số của dao động là $\nguon{ly11 kntt.pdf}$
\choice
{\True $5 \mathrm{Hz}$.}
{$10 \mathrm{Hz}$.}
{$20 \mathrm{Hz}$.}
{5 $\pi$ Hz.}
\loigiai{
Tần số f = $\dfrac{\omega}{2\pi}=5\,\mathrm{Hz}$
}
\end{ex}

% ===== Câu 18 =====
% ID: L11C1B1-183-DS
% Mức: TH
\begin{ex}
% ID: L11C1B1-183-DS
% Mức: TH
Một chất điểm dao động điều hoà với chu ki $T = 2\$, $\mathrm{s}$. Trong $3\,\mathrm{s}$ vật đi được quãng đường $60 \mathrm{cm}$. Khi $t = 0 vật đi qua vị tri cân bằng và$ hướng về vị trí biên dương. Hãy viết phương trình dao động của vật. 5 $\pi\nguon{ly11 kntt.pdf}$
\choiceTF

\loigiai{
Có $T=2\,\mathrm{s}; t =3\,\mathrm{s}= 1,5T$, do vậy quãng đường vật dao động đi được sau t = $3\,\mathrm{s}$ là:

s = $4\,\mathrm{A}$ + $2\,\mathrm{A}$ = $6\,\mathrm{A}$ = $60\,\mathrm{cm}\RightarrowA$ = $10\,\mathrm{cm}$; $\omega = \dfrac{2\pi}{T} = \dfrac{2\pi}{2} = \pi\left( \text{rad/s} \right);$ khi t = 0 vật qua vị trí cân bằng và đang hướng về phía biên dương, nên $\varphi = - \dfrac{\pi}{2}$.

Suy ra phương trình dao động của vật là: x = 10cos $\left( {\pi t - \dfrac{\pi}{2}} \right)$ (cm).
}
\end{ex}

% ===== Câu 20 =====
% ID: L11C1B1-184-TN
% Mức: NB
\begin{ex}
% ID: L11C1B1-184-TN
% Mức: NB
Một chất điểm dao động điều hoà. Trong thời gian 1 phút, vật thực hiện được 30 dao động. Chu kì dao động của chất điểm là $\nguon{ly11 kntt.pdf}$
\choice
{\True $2\,\mathrm{s}$.}
{$30\,\mathrm{s}$.}
{$0,5\,\mathrm{s}$.}
{$1\,\mathrm{s}$.}
\loigiai{
Chu kì $T=\dfrac{t}{N} = \dfrac{60}{30}=2\,\mathrm{s}$
}
\end{ex}

% ===== Câu 22 =====
% ID: L11C1B1-185-DS
% Mức: TH
\begin{ex}
% ID: L11C1B1-185-DS
% Mức: TH
Một chất điểm dao động điều hoà theo phương trình $x=10\cos\left(2\pi t+\dfrac{5\pi}{6}\right)\,(\mathrm{cm})$. Tính quãng đường vật đi được trong khoảng thời gian từ $t_1=1\,\mathrm{s}$ đến $t_2=2,5\,\mathrm{s}$. $\nguon{ly11 kntt.pdf}$
\choiceTF

\loigiai{
$T = \dfrac{2\pi}{\omega} = \dfrac{2\pi}{2\pi}=1\,\mathrm{s}\Rightarrowt₁ =1\,\mathrm{s}$ = 1 T.

Quãng đường vật đi được sau t₁ = $1\,\mathrm{s}$ là: s₁ = $4\,\mathrm{A}$ = $40\,\mathrm{cm}$ t₂ = $2,5\,\mathrm{s}$ = 2,5 $T$ = 2 $T$ + $\dfrac{T}{2}$ Quãng đường vật đi được sau t₂ = 2, $5\,\mathrm{s}$ là: s₁ = 2. $4\,\mathrm{A}$ + $2\,\mathrm{A}$ = $10\,\mathrm{A}$ = $100\,\mathrm{cm}$

Quãng đường vật đi được từ t₁ đến t₂ là: $\Delta$ s = s₂ - s₁ = 100 - 40 = $60\,\mathrm{cm}$.
}
\end{ex}

% ===== Câu 24 =====
% ID: L11C1B1-186-TN
% Mức: TH
\begin{ex}
% ID: L11C1B1-186-TN
% Mức: TH
Một chất điểm dao động điều hoà có phương trình li độ theo thời gian là $x=5\sqrt{3}\cos\left(10\pi t+\dfrac{\pi}{3}\right)\,(\mathrm{cm})$. Tần số của dao động là $\nguon{ly11 kntt.pdf}$
\choice
{$10 \mathrm{Hz}$.}
{$20 \mathrm{Hz}$.}
{10 $\pi$ Hz.}
{\True $5 \mathrm{Hz}$.}
\loigiai{
Phương trình dao động điều hòa có dạng $x = A\cos(\omega t + \phi)$. Từ phương trình đã cho $x=5\sqrt{3}\cos\left(10\pi t+\dfrac{\pi}{3}\right)\,(\mathrm{cm})$, ta xác định được tần số góc $\omega = 10\pi \, \mathrm{rad/s}$. Tần số $f$ của dao động được tính bằng công thức $f = \dfrac{\omega}{2\pi}$. Thay số vào ta có $f = \dfrac{10\pi}{2\pi} = 5 \, \mathrm{Hz}$.
}
\end{ex}

% ===== Câu 28 =====
% ID: L11C1B1-187-TN
% Mức: TH
\begin{ex}
% ID: L11C1B1-187-TN
% Mức: TH
Một chất điểm dao động điều hoà có phương trình li độ theo thời gian là $x=6\cos\left(4\pi t+\dfrac{\pi}{3}\right)\,(\mathrm{cm})$. Chu kì của dao động là $\nguon{ly11 kntt.pdf}$
\choice
{$4\,\mathrm{s}$.}
{$2\,\mathrm{s}$.}
{$0,25 \mathrm{cm}$.}
{\True $0,5\,\mathrm{s}$.}
\loigiai{
Chọn D.

Từ phương trình xác định được tần số góc $. \omega = 4\pi rad/s\Rightarrow T = \dfrac{2\pi}{\omega} = 0,$ 5 $\mathrm{s}$
}
\end{ex}

% ===== Câu 32 =====
% ID: L11C1B1-188-TN
% Mức: TH
\begin{ex}
% ID: L11C1B1-188-TN
% Mức: TH
Một chất điểm dao động điều hoà có phương trình li độ theo thời gian là $x=10\cos\left(\dfrac{\pi}{3}t+\dfrac{\pi}{2}\right)\,(\mathrm{cm})$. Tại thời điểm $t$, vật có li độ $6\,\mathrm{cm}$ và đang hướng về vị trí cân bằng. Sau $9\,\mathrm{s}$ kể từ thời điểm $t$ thì vật đi qua li độ $\nguon{ly11 kntt.pdf}$
\choice
{$3 \mathrm{cm}$ đang hướng về vị trí cân bằng}
{$-3 \mathrm{cm}$ đang hướng về vị trí biên.}
{$6 \mathrm{cm}$ đang hướng về vị trí biên.}
{\True $-6 \mathrm{cm}$ đang hướng về vị trí cân bằng. $\pi$}
\loigiai{
Chọn D.

Chu kì $T=\dfrac{2\pi}{\omega}=6\,\mathrm{s}\Rightarrowt₁ =9\,\mathrm{s}= 1,5T.

Sau t₁ = 1,5T$vật ở vị trí như Hình$2.1\,\text{G}$.

[Một chất điểm dao động điều hoà có phương trình li độ theo thời gian là x = 10cos $\pi /3t+\pi /2)$]

Do đó sau $9\,\mathrm{s}$ kể từ thời điểm t thì vật đi qua li độ là - $6\,\mathrm{cm}$ và đang hướng về vị trí cân bằng.
}
\end{ex}

% ===== Câu 2 =====
% ID: L11C1B1-189-DS
% Mức: TH
\begin{ex}
% ID: L11C1B1-189-DS
% Mức: TH
Phương trình dao động điều hoà là $x=5\cos(2\pi t+\dfrac{\pi}{3}
ight)\,(\mathrm{cm})$. Hãy cho biết biên độ, pha ban đầu và pha ở thời điểm $t$ của dao động. $\nguon{ly11 kntt.pdf}$
\choiceTF
{\True Biên độ là $5\,\mathrm{cm}$.}
{\True Pha ban đầu là $\dfrac{\pi}{3}\,\mathrm{rad}$.}
{\True Pha ở thời điểm $t$ là $2\pi t + \dfrac{\pi}{3}\,\mathrm{rad}$.}
{\True Tất cả đều đúng.}
\loigiai{
\begin{itemchoice}
\item (Đúng) Biên độ là $5\,\mathrm{cm}$. (Vì): Biên độ $A = 5\,\mathrm{cm}$. Mệnh đề $A$ đúng
\item (Đúng) Pha ban đầu là $\dfrac{\pi}{3}\,\mathrm{rad}$. (Vì): Pha ban đầu $\varphi = \dfrac{\pi}{3}\,\mathrm{rad}$. Mệnh đề $B$ đúng
\item (Đúng) Pha ở thời điểm $t$ là $2\pi t + \dfrac{\pi}{3}\,\mathrm{rad}$. (Vì): Pha ở thời điểm $t$ là $\omega t + \varphi = 2\pi t + \dfrac{\pi}{3}\,\mathrm{rad}$. Mệnh đề $C$ đúng
\item (Đúng) Tất cả đều đúng. (Vì): Vì cả ba mệnh đề $A$, $B$, $C$ đều đúng nên mệnh đề $D$ đúng
\end{itemchoice}
}
\end{ex}

% ===== Câu 4 =====
% ID: L11C1B1-190-TN
% Mức: TH
\begin{ex}
% ID: L11C1B1-190-TN
% Mức: TH
Một chất điểm dao động điều hoà với phương trình $x=5\cos\left(10\pi t+\dfrac{\pi}{3}\right)\,(\mathrm{cm})$. Li độ của chất điểm khi pha dao động bằng $\pi$ là $\nguon{ly11 kntt.pdf}$
\choice
{$5 \mathrm{cm}$.}
{\True $-5 \mathrm{cm}$.}
{$2,5 \mathrm{cm}$.}
{$-2,5 \mathrm{cm}$.}
\loigiai{
Pha dao động là $\pi$. Li độ của chất điểm được tính bằng công thức $x = A\cos(\phi)$. Thay số vào công thức, ta có $x = 5\cos(\pi) = 5 \times (-1) = -5 \, (\mathrm{cm})$.
Li độ của chất điểm khi pha dao động bằng $\pi$ là $-5 \mathrm{cm}$
}
\end{ex}

% ===== Câu 5 =====
% ID: L11C1B1-191-DS
% Mức: TH
\begin{ex}
% ID: L11C1B1-191-DS
% Mức: TH
Một chất điểm dao động điều hoà có phương trình li độ theo thời gian là $x=10\cos\left(\dfrac{\pi}{3}t+\dfrac{\pi}{2}\right)\,(\mathrm{cm})$. a) Tính quãng đường chất điểm đi được sau 2 dao động. b) Tính li độ của chất điểm khi $t=6\,\mathrm{s}$. $\nguon{ly11 kntt.pdf}$
\choiceTF

\loigiai{
a) Quãng đường vật đi được sau 2 dao động: s = 2. $4\,\mathrm{A}= 2.4.10 =80\,\mathrm{cm}.

b) Khi t =6\,\mathrm{s}\Rightarrow$x = 10cos$\left( {2\pi.6 + \dfrac{\pi}{2}} \right)$ = 0.
}
\end{ex}

% ===== Câu 7 =====
% ID: L11C1B1-192-TN
% Mức: TH
\begin{ex}
% ID: L11C1B1-192-TN
% Mức: TH
Một chất điểm dao động điều hoà có phương trình li độ theo thời gian là $x=5\sqrt{3}\cos\left(10\pi t+\dfrac{\pi}{3}\right)\,(\mathrm{cm})$. Tại thời điểm $t=1\,\mathrm{s}$ thì li độ của chất điểm bằng $\nguon{ly11 kntt.pdf}$
\choice
{$2,5 \mathrm{cm}$.}
{$-5\sqrt{3} \mathrm{cm}$.}
{$5 \mathrm{cm}$.}
{\True $2,5\sqrt{3} \mathrm{cm}$.}
\loigiai{
Chọn D.

Tại thời điểm t = $1\,\mathrm{s}$ thì li độ của chất điểm bằng x = 5 $\sqrt{3}$ cos $\left( {10\pi.1 + \dfrac{\pi}{3}} \right)$ = 2,5 $\sqrt{3}$ (cm)
}
\end{ex}

% ===== Câu 10 =====
% ID: L11C1B1-193-TN
% Mức: TH
\begin{ex}
% ID: L11C1B1-193-TN
% Mức: TH
Một chất điểm dao động điều hoà có phương trình li độ theo thời gian là $x=6\cos\left(10\pi t+\dfrac{\pi}{3}\right)\,(\mathrm{cm})$. Li độ của chất điểm khi pha dao động bằng $-\dfrac{\pi}{3}$ là $\nguon{ly11 kntt.pdf}$
\choice
{\True $3 \mathrm{cm}$.}
{$-3 \mathrm{cm}$.}
{$3\sqrt{3} \mathrm{cm}$.}
{$-3\sqrt{3} \mathrm{cm}$.}
\loigiai{
Chọn A.

Li độ của chất điểm khi pha dao động bằng $\left( {- \dfrac{\pi}{3}} \right)$ là x = 6cos $\left( {- \dfrac{\pi}{3}} \right)$ = 3 (cm)
}
\end{ex}

% ===== Câu 13 =====
% ID: L11C1B1-194-TN
% Mức: TH
\begin{ex}
% ID: L11C1B1-194-TN
% Mức: TH
Phương trình dao động của một vật có dạng $x=-A\cos\left(\omega t+\dfrac{\pi}{3}\right)\,(\mathrm{cm})$. Pha ban đầu của dao động là $\nguon{ly11 kntt.pdf}$
\choice
{$\dfrac{\pi}{3}$ rad}
{$-\dfrac{\pi}{3}$ rad}
{$\dfrac{2\pi}{3}$ rad}
{\True $-\dfrac{2\pi}{3}$ rad}
\loigiai{
Chọn D.

$x = - A\cos\left( {\omega t + \dfrac{\pi}{3}} \right) = A\cos\left( {\pi - \omega t - \dfrac{\pi}{3}} \right) = A\cos\left( {\omega t - \dfrac{2\pi}{3}} \right)\left( \text{cm} \right).$

Pha ban đầu là - $\dfrac{2\pi}{3}$ rad.
}
\end{ex}
